% Generated draft: requires mathematical and source review.
\documentclass[11pt]{article}
\usepackage[margin=25mm]{geometry}
\usepackage{fontspec}
\setmainfont{Latin Modern Roman}
\usepackage{amsmath,amssymb,mathtools,unicode-math}
\setmathfont{Latin Modern Math}
\usepackage{enumitem,needspace,xurl,hyperref}
\setlength{\parindent}{0pt}
\setlength{\parskip}{5pt}
\setlength{\emergencystretch}{3em}
\newcommand{\N}{\mathbb{N}}
\newcommand{\Z}{\mathbb{Z}}
\newcommand{\Q}{\mathbb{Q}}
\newcommand{\R}{\mathbb{R}}
\newcommand{\C}{\mathbb{C}}
\newcommand{\F}{\mathbb{F}}
\newcommand{\A}{\mathbb{A}}
\newcommand{\PP}{\mathbb{P}}
\newcommand{\E}{\mathbb{E}}
\newcommand{\eps}{\varepsilon}
\newcommand{\dd}{\,\mathrm{d}}
\begin{document}
\section*{UnsolvedMath ID: 30005390}
\textbf{Source problem code:} catalogue entry
\subsection{Definitions and mathematical statement}
Let \(\ell\) be a prime and let \(\F_\ell\) denote the field with \(\ell\) elements. The source uses spaces \(X_K\) and \(X_{K_0(\mathfrak p)}\), a cohomological degree \(d\), a Hecke ideal \(\mathfrak m\), and two maps denoted \(\pi_{1*}\) and \(\pi_{2*}\). Their full geometric and Hecke-theoretic setup is absent from the supplied excerpt. In particular, the number field, group, allowed levels, prime \(\mathfrak p\), coefficient restrictions, and meaning of non-Eisenstein are not specified sufficiently to give an unconditional standalone formulation.

For a commutative Hecke algebra \(\mathbb T\), a maximal ideal \(\mathfrak m\subset\mathbb T\), and a \(\mathbb T\)-module \(M\), write \(M_{\mathfrak m}=(\mathbb T\setminus\mathfrak m)^{-1}M\). This explains the localization notation; the intended algebra and its compatible actions on the two cohomology groups still require specification. The surrounding source discussion refers to mod \(\ell\) \'etale cohomology, but does not supply a complete definition of the cohomology theory or geometric objects in the displayed conjecture.

\textbf{Conjecture 2 (generalized Ihara lemma, as supplied).} In the setting intended by the source, if \(\mathfrak m\) is non-Eisenstein, then
\[
\Phi_{\mathfrak m}:=\pi_{1*}+\pi_{2*}:
H^d(X_K,\F_\ell)_{\mathfrak m}^{\oplus 2}
\longrightarrow
H^d(X_{K_0(\mathfrak p)},\F_\ell)_{\mathfrak m},
\qquad
(a,b)\longmapsto\pi_{1*}(a)+\pi_{2*}(b),
\]
is injective. Explicitly, for every \(a,b\in H^d(X_K,\F_\ell)_{\mathfrak m}\),
\[
\pi_{1*}(a)+\pi_{2*}(b)=0
\quad\Longrightarrow\quad a=b=0.
\]
The subscripts \(*\) and the direction of the displayed map are retained exactly from the supplied statement. The excerpt does not define the underlying degeneracy morphisms or explain whether these symbols denote pullbacks, pushforwards, or another induced operation; this must be resolved before the map is fully specified. The admissible values of \(d\), its relationship to the dimension, and any restrictions involving \(\ell\) and \(\mathfrak p\) are also unresolved.

\textbf{Supplied qualifications and related formulations.} The excerpt reports that the \(d=0\) case follows from strong approximation; that Manning and Shotton proved the \(d=1\) case under an unspecified mild technical assumption; and that arbitrary-dimensional results are known roughly under a Fontaine--Laffaille condition on mod \(\ell\) \'etale cohomology, which can impose strong restrictions on \(\ell\). These are source claims, not hypotheses to be silently added to the conjecture.

The source also describes a direct limit, schematically
\[
V^d_{\mathfrak m}
=\varinjlim_{K'_{\mathfrak p}}
H^d(X_{K^{\mathfrak p}K'_{\mathfrak p}},\F_\ell)_{\mathfrak m},
\]
as an admissible smooth mod \(\ell\) representation of \(\operatorname{GL}_2(F_{\mathfrak p})\). Here the source appears to hold the level away from \(\mathfrak p\) fixed and vary the local level; the indexing groups and transition maps are not supplied. The direct-limit notation above is a schematic transcription of a corrupted formula, not a fully defined additional object. According to the excerpt, Ihara's lemma is equivalent, at least after assuming its unspecified Conjecture 1, to every irreducible subrepresentation of this representation being generic when \(\mathfrak m\) is non-Eisenstein. The local field, coefficient conventions, and definition of genericity require the missing setup.

In a separate remark assuming \(\ell\ne p\), \(F=\Q\), and \(G=\operatorname{GL}_2\), the source relates Ihara's lemma together with a mod \(\ell\) multiplicity-one result to Emerton's compatibility of \(\ell\)-adically completed cohomology with the local Langlands correspondence in families of Emerton and Helm at \(p\). It also describes the representation-theoretic formulation as important for compatibility with categorical local Langlands in work of Emerton--Gee--Zhu described there as upcoming. The source explicitly allows nongeneric, described there as one-dimensional, representations as subquotients of its cohomological representation; this does not replace the assertion about irreducible subrepresentations. The precise cohomological notation in that sentence is corrupted.

The excerpt additionally states a separate Theorem 1: if \(\ell\) is odd and the image of \(\bar\rho_{\mathfrak m}\) is nonsolvable, then Conjecture 1 holds. Neither \(\bar\rho_{\mathfrak m}\) nor Conjecture 1 is defined in the supplied record. This assertion therefore cannot be substituted for the conjecture above. The following sentence is truncated.
\subsection{Short English statement}
The proposed generalized Ihara lemma says that, at a non-Eisenstein Hecke ideal, the sum of the two indicated maps from two copies of level-\(K\) cohomology to level-\(K_0(\mathfrak p)\) cohomology is injective. The supplied excerpt omits essential setup and map conventions, so this is a statement-only transcription with unresolved qualifications requiring human mathematical review. The registry labels the problem partially solved; that label and its reported special cases are supplied assessments.
\subsection{Sources}
\begin{itemize}
\item Supplied local UnsolvedMath registry record, permanent ID \(30005390\), title \emph{Generalized Ihara Lemma for Locally Symmetric Spaces}, exact source problem code \texttt{catalogue entry}; collection: Oberwolfach Reports Open Problems; proposed year: 2023. No individual proposer is identified.
\item Supplied report citation: \emph{Arithmetic of Shimura Varieties} (2023), Oberwolfach Reports, DOI 10.4171/owr/2023/5, pp.~261--326. \url{https://doi.org/10.4171/owr/2023/5}. The excerpt identifies the target as Conjecture 2. The record labels its statement \texttt{corrected\_verified} and says it was reconstructed and checked against the official report; this is supplied verification metadata, not independent verification here.
\item Supplied literature reference, labeled ``Manning and Shotton, Ihara lemma for Hilbert modular varieties'': \url{https://arxiv.org/abs/math/0503134}. The correspondence between this label, this URL, and the excerpt's reference [6] requires verification; no bibliographic correction is inferred.
\item Supplied literature reference, labeled ``Oberwolfach Report 5/2023'': \url{https://doi.org/10.4171/owr/2023/5}.
\item Literature assessment dated exactly \texttt{2026-08-22}: the supplied curation labels the problem partially solved and reports results in dimension zero, dimension one under hypotheses, and several higher-dimensional Fontaine--Laffaille or large-image cases, but not in the full stated generality. This is supplied dated curation, not independently verified; the record itself requests independent verification of time-sensitive status and cited claims.
\item Further attribution in the supplied excerpt: applications to level-raising are traced to Ribet, and applications to congruences between special values of \(L\)-functions are mentioned. Clozel--Harris--Taylor [4] are credited with using a representation-theoretic formulation to propose higher-rank generalizations for non-minimal modularity lifting. The remark attributes local Langlands-in-families applications to Emerton and to Emerton--Helm [5], and mentions Emerton--Gee--Zhu work as upcoming at the time of the source. The separate Theorem 1 is attributed to [2], using a modification of a unitary Shimura variety strategy from [1]. Full bibliographic entries for [1], [2], [4], [5], and [6] are absent from the supplied record.
\end{itemize}
\end{document}
